\section{Introduction}
\label{sec:intro}
Recent advances in video understanding~\cite{reid2024gemini, liu2025st, hong2024cogvlm2, liu2024oryx, zhang2024STGT, li2024aria} and generation~\cite{SORA, vidu, kelin, hunyuanvideo} have been driven by large vision-language models (VLMs). For video comprehension, researchers have extended image-based architectures (\eg, LLaVA-Video~\cite{zhang2024video}, PLLaVA~\cite{xu2024pllava} and CogVLM2-Video~\cite{hong2024cogvlm2}) and explored hybrid image-video training approaches (Qwen2-VL~\cite{wang2024qwen2}, LLaVA-OneVision~\cite{li2024llava}). Currently, video generation systems such as Sora~\cite{SORA} and HunyuanVideo~\cite{hunyuanvideo} use VLMs for multimodal captioning and prompt engineering. However, current benchmarks~\cite{li2024mvbench, li2024videovista, wang2024lvbench, fu2024video, zhou2024mlvu, wu2024longvideobench} struggle to assess the detailed spatio-temporal aspects required for these applications.


In visual generation, existing caption evaluation metrics fall mainly into two categories: reference-based (METEOR~\cite{meteor}, BLEU~\cite{bleu}, SPICE~\cite{Spice}, and CIDEr~\cite{cider}) and reference-free (InfoMetIC~\cite{infometic}, CLIPScore~\cite{Clipscore}, and TIGEr~\cite{Tiger}). Reference-based metrics~\cite{meteor,bleu,Spice,cider} assess the quality of the captions by comparing them with ground-truth captions. However, these scores are highly dependent on the reference captions' format. Hence, FAIEr~\cite{Faier} adopts visual and textual scene graphs for a more robust reference-based evaluation. Reference-free metrics~\cite{infometic,Clipscore,Tiger} use semantic vectors from the reference image to evaluate the similarity of the caption. These methods falter with concept-dense captions, overwhelmed by numerous concepts. Recent innovations like QACE~\cite{QACE}, DSG ~\cite{cho2023davidsonian}], DPG-bench~\cite{Ella} and CapsBench~\cite{playgroundv3} employ question-answering frameworks to address dense concept evaluation for image captioning. Notably, these methods derived from image description evaluation still have significant measurement blind spots when dealing with the spatio-temporal dynamic elements of video generation scenarios, \eg,  camera motion (``slow zoom" vs. ``fast pan") or dynamic spatial relationships (object displacement from ``left foreground " to ``central midground"). 
In these cases, video generation systems (\eg, Sora) need to ensure accurate text-video alignment.

Current video understanding benchmarks primarily focus on holistic semantic alignment (MVBench ~\cite{li2024mvbench}, VideoVista~\cite{li2024videovista}) or specific skill evaluation (LVBench~\cite{wang2024lvbench}), neglecting fine-grained spatial-temporal dynamics essential for video generation. This gap becomes critical when evaluating text-to-video systems where caption precision directly impacts visual output. For instance, failing to distinguish between ``slow zoom" versus ``fast pan" camera motions, or misrepresenting object trajectories from ``left foreground" to ``central midground". While manual evaluation of such nuances remains impractical at scale, automated metrics also struggle with dynamic semantic alignment due to inherent limitations in traditional evaluation paradigms. To address this problem, it is necessary to re-examine the design principles of video description evaluation metrics.

To address critical gaps in evaluating video caption,  we introduce \textbf{VCapsBench}, the first large-scale fine-grained benchmark for video caption evaluation. It contains 5,677 diverse videos with 109,796 human-verified questions, allowing fine-grained evaluation of caption quality. Specifically, we employ text-based question-answering across 21 categories, including action, camera movements (\eg, zoom, pan and tilt), object positioning (absolute or relative position), entity and shot type, \etc, with ``yes", ``no", and ``unanswerable" ternary judgments. Unlike image-focused caption benchmarks like CapsBench~\cite{playgroundv3}, our VCapsBench prioritizes temporal continuity through video-specific queries (avg. 19 questions/video) while mitigating LLM hallucination via an ``Unanswerable" option. To accurately and comprehensively assess caption quality, we introduce three metrics: Accuracy (AR), Inconsistency Rate (IR), and Coverage Rate (CR). To automate the calculation of these metrics, we develop an evaluation pipeline that leverages powerful LLMs for objective video caption quality measurement, as illustrated in Fig.~\ref{fig:fig1}. In summary, our contributions are as follows:
\begin{itemize}
    \item We introduce \textbf{VCapsBench}, the first large-scale fine-grained benchmark for video caption evaluation, featuring diverse videos (5K+) and QA-pairs (100K+).
    \item We introduce three metrics: Accuracy (\textbf{AR}), Inconsistency Rate (\textbf{IR}), and Coverage Rate (\textbf{CR}), along with an automated evaluation pipeline to fairly assess the correctness and coverage of video captions.
    
    \item We evaluate ten VLMs, including seven advanced open models (Qwen2VL, Qwen2.5VL, InternVL2.5, LLaVA-Video, and VideoLLaMA3, \etc) and three closed model, GPT-4o,  Gemini2.5-Pro Flash and Preview, providing a solid reference for the community.
\end{itemize}





